\scanpage{4}
\mspage{3}
Alors on a, \deleted{si $X$ purement de dim $n$ et $i\leq n$ (resp. $i\geq n$)}
\[
\begin{aligned}
\mathrm{Filt}^{\unsure{j|d}} H^i(\bar X,Q_\ell)&=\text{partie de }T^\ell(h^i(X))\ \text{\deleted{provenant}\inserted{image} du}\\
&\quad\text{plus grand sous-motif \inserted{de}}\\
&\quad\text{niveau }\leq i-2j\ \text{\deleted{(resp. $2n-i\unsure{-j|\geq j}$)} \deleted{de} de }h^i(X)
\end{aligned}
\]
Cette caractérisation \inserted{du niveau des motifs} est tout au moins correcte si on admet la résolution des singularités, \deleted{\unsure{sur $k$}} \inserted{des variétés de dim $\leq\dim X=n$ sur $k$.} Comme conséquence de la résolution et des conjectures standard, on trouve donc que les dimensions des $\mathrm{Filt}^j H^i(\bar X,Q_\ell)$ sont indépendantes de $\ell$. (-- Bien mieux, ils correspondent à des motifs, et les quotients facteurs successifs sont \deleted{\unsure{et}} purs de poids $i-2j$ resp. $2n-i-j$).

\inserted{Parler des inv. birationnels\ldots}
% slanted
\margin{\uuline{Conjectures de Tate et de Hodge pourraient passer ici}}
7. \uline{Structure des} $\otimes$-\uline{catégories}. \deleted{\unsure{abéliennes|standard}}

\deleted{$\otimes$ catégorie.} Catégorie abélienne avec opération $\otimes$ ``ass'', ``comm'' et ``unitaire''. On suppose de plus l'existence de \uline{Hom} internes, au sens du $\otimes$.

On dit qu'elle est ``spéciale'' si
\[
\check M\otimes N\xrightarrow{\sim}\underline{\mathrm{Hom}}(M,N)
\]
Supposons de plus que $C$ soit $K$-linéaire ($K$ un corps) (\unsure{i.e.} on se donne $K\longrightarrow\operatorname{End}(\mathrm{id}_{C})$), et que les Hom soient de dim finie sur $k$ : \deleted{\unsure{on est}} On dira qu'on a une $\otimes$-catégorie sur $k$ (pour simplifier).

Une telle catégorie est \unsure{dite|dit} ``connexe'' si $\operatorname{End}(\mathbf{1}_C)\simeq k$ (\unsure{par|bien} l'identité).
% archive sheet number: 3

